\chapter{Conclusions}
\label{ch:conclusions}

The present work defines a CFD-based framework for investigating the influence of gas properties, inlet velocity and obturator-hole geometry on excess flow valve behavior. The reference Design of Experiments is the \texttt{SIM\_DoE\_2210} campaign, comprising 24 planned combinations of two gases, three inlet velocities ($5$, $10$ and $15\,\mathrm{m\,s^{-1}}$), two hole counts and two hole diameters.

The study combines OpenFOAM 10 simulations with a scripted preparation and HPC workflow. The intended response variables include the inlet-to-outlet pressure difference, the axial fluid load acting on a stationary obturator and the corresponding dimensionless coefficients.

\paragraph{Status of the numerical investigation.}
At the time of drafting, a subset of the hydrogen cases still requires numerical stabilization. The reported forces also require confirmation of the density used by the incompressible force integration. Consequently, final statements about geometric ranking, cross-fluid performance, pressure-loss trends or valve trip behavior are deferred until convergence, conservation, mesh quality and density conventions have been checked. Reaching the final prescribed solver iteration is not sufficient evidence of convergence.

\paragraph{Scope and future developments.}
The fixed-geometry, incompressible steady-state approach is intended to support a comparative design assessment, not to predict the transient closing trajectory or closing time of a complete mechanical valve. After verification of the dataset, future work may complete the DoE response analysis and the interactive comparison tool, evaluate locally compressible flow effects and investigate transient or fluid--structure interaction methods to resolve obturator motion.

% TODO: Replace this provisional conclusion with quantified, verified findings
% after the SIM_DoE_2210 campaign and the interactive tool are completed.
